\chapter{Introdução aos Sistemas Embarcados}
\label{cap:embarcados}

Com os fundamentos da linguagem C estabelecidos, a Parte II desta apostila muda de ambiente: em vez de compilar e rodar programas em um computador comum, vamos escrever firmware para um microcontrolador real, o \ESP{}. Este capítulo estabelece o vocabulário e o contexto necessários antes de abrirmos o PlatformIO pela primeira vez, no próximo capítulo.

\section{O que é um sistema embarcado}

Um \textbf{sistema embarcado} é um sistema computacional dedicado a uma função específica, geralmente integrado a um equipamento maior — diferente de um computador de propósito geral, que roda uma variedade praticamente ilimitada de programas. Uma máquina de lavar, um roteador Wi-Fi, o painel de instrumentos de um carro, uma pulseira de atividade física: todos são sistemas embarcados.

\begin{table}[h]
\centering
\begin{tabular}{@{}lll@{}}
\toprule
\textbf{Característica} & \textbf{Computador comum} & \textbf{Sistema embarcado} \\
\midrule
Propósito       & genérico            & dedicado a uma função \\
Memória RAM     & gigabytes           & kilobytes a poucos megabytes \\
Armazenamento   & discos de alta capacidade & memória flash, tipicamente MB \\
Sistema operacional & completo (Linux, Windows) & ausente, mínimo, ou um RTOS \\
Entrada/saída   & teclado, mouse, tela & sensores, atuadores, LEDs \\
Consumo de energia & watts        & miliwatts a microwatts \\
\bottomrule
\end{tabular}
\caption{Sistemas embarcados vs. computadores de propósito geral.}
\end{table}

Essas restrições — pouca memória, sem sistema operacional completo, energia limitada — são exatamente o motivo pelo qual C continua sendo a linguagem dominante nesse domínio: ela permite controle fino sobre o que o processador de fato executa, sem a sobrecarga (\emph{overhead}) de linguagens de mais alto nível.

\section{O microcontrolador ESP32}

O \ESP{} é uma família de \textbf{microcontroladores} (\emph{system-on-a-chip}, SoC) desenvolvida pela Espressif Systems, que integra em um único chip:

\begin{itemize}
    \item Um processador \textbf{Xtensa LX6 (ou LX7, em variantes mais recentes)}, tipicamente dual-core, operando a até 240~MHz;
    \item \textbf{Wi-Fi} (802.11 b/g/n) e \textbf{Bluetooth} (clássico e BLE) integrados;
    \item Memória \textbf{SRAM} interna (algumas centenas de kB) para variáveis e pilha durante a execução;
    \item Memória \textbf{Flash} externa (tipicamente 4~MB ou mais) onde o firmware compilado é armazenado permanentemente;
    \item Dezenas de pinos \textbf{GPIO} de uso geral, além de periféricos dedicados: conversores analógico-digitais (ADC), saídas PWM, interfaces I²C, SPI, UART, e mais.
\end{itemize}

\begin{nota}
``Microcontrolador'' e ``microprocessador'' não são sinônimos. Um \textbf{microprocessador} (como o de um computador) é apenas a unidade de processamento — precisa de chips externos para memória, entrada/saída, etc. Um \textbf{microcontrolador} já integra processador, memória e periféricos em um único chip, pronto para ser incorporado a um produto com o mínimo de componentes adicionais. É exatamente esse nível de integração que torna o \ESP{} tão popular em projetos de eletrônica e produtos comerciais.
\end{nota}

\section{Como um firmware é diferente de um programa comum}

Ao escrever \code{hello\_world.c} no \cref{cap:primeiros-passos}, o programa era executado sob um sistema operacional (Linux, Windows ou macOS), que gerenciava memória, agendamento de processos e acesso a dispositivos por trás dos panos. Em um microcontrolador programado da forma mais simples, não existe esse intermediário: o código compilado é o \textbf{único} programa em execução, com acesso direto ao hardware.

Isso muda a estrutura típica de um programa de duas formas centrais:

\begin{enumerate}
    \item Não existe uma função \code{main} que ``termina'' — como adiantado no \cref{cap:primeiros-passos}, o firmware gira indefinidamente em torno de um \textbf{laço principal}, reagindo a eventos e atualizando o estado do sistema enquanto houver energia;
    \item O programador é responsável por configurar diretamente os periféricos que deseja usar — não existe um driver de teclado ou de tela pronto: existe, por exemplo, um pino GPIO específico que precisa ser configurado como saída antes de poder acender um LED.
\end{enumerate}

\begin{esp32box}
Tecnicamente, o \ESP{} \emph{não} roda ``sem sistema operacional'' — por baixo do framework que usaremos nesta apostila, o \textbf{Arduino}, existe o \textbf{ESP-IDF}, o SDK oficial da Espressif, construído sobre o \textbf{FreeRTOS}, um sistema operacional de tempo real (RTOS) leve, projetado especificamente para microcontroladores. O framework Arduino é, na prática, uma camada de conveniência sobre essa base: ele simplifica drasticamente a forma de escrever o firmware — como veremos já no próximo capítulo — sem esconder o acesso direto ao hardware que torna C tão adequado para essa tarefa.
\end{esp32box}

Antes de colocar a mão no código, o próximo capítulo faz uma pausa conceitual: de onde vêm os microcontroladores, e como sua unidade de processamento funciona por dentro. Esse pano de fundo prepara o terreno para o \textbf{PlatformIO}, o ambiente que usaremos a partir do \cref{cap:platformio} para escrever, compilar e carregar código C no \ESP{}.
